% 1. 注册命名空间
\installnamespace{ScHead}

% 2. 安装处理器
\installcommandhandler \XXXXMyNote {MyNote} \XXXXMyNote

% 3. 设置默认参数
\setupMyNote
  [title=笔记,
   color=blue,
   width=\textwidth]


%%%%%%%%%%%%%%%%%%%%%%%%%%%%%%%%%%%%%%%%%%%%%%%%%%%%%%%%%%%%%%%%%%%%%%%%%%%%%%%%%%

\unprotect

% 定义 \startMyNote
\def\startMyNote[#1]{
    \begingroup % 开启局部作用域，确保参数修改不影响环境外
    % 检查 #1 是不是 key=value 赋值；如果是，则临时应用这些设置
    \doifassignmentelse{#1}{\setupMyNote[#1]}{\def\currentMyNote{#1}}

    % 这里调用 ConTeXt 的排版组件来承载内容
    \startframedtext[
        title=\ScHeadparameter{title},
        foregroundcolor=\ScHeadparameter{color},
        framecolor=\ScHeadparameter{color},
        width=\ScHeadparameter{width}
    ]
}

% 定义 \stopMyNote
\def\stopMyNote{
    \stopframedtext
    \endgroup % 结束局部作用域
}

\protect



% usage:
\starttext

% 1. 使用默认值
\startMyNote
  这段文字使用默认的蓝色和标题。
\stopMyNote

\blank

% 2. 使用命名参数自定义（即时生效）
\startMyNote[title={重要提示}, color=red, width=10cm]
  这段文字变成了红色，标题也变了，且宽度变窄。
\stopMyNote

\stoptext